\documentclass[11pt]{article}
\usepackage[margin=2.2cm]{geometry}
\usepackage[T1]{fontenc}
\usepackage{amsmath,amssymb,booktabs,xcolor,enumitem,caption,array}
\usepackage[section]{placeins}
\usepackage{hyperref}
\usepackage{xurl}
\usepackage{fvextra}
\fvset{breaklines=true,fontsize=\footnotesize}
\setlist{itemsep=2pt,topsep=3pt}
\hypersetup{colorlinks=true,linkcolor=blue!50!black,urlcolor=blue!50!black}
\captionsetup{font=small,labelfont=bf}
\setlength{\emergencystretch}{3em}
\newcommand{\se}{\sigma_e}

\title{Machine learning for CCD commissioning: a plan, and how to test it on the data we have}
\author{ccd\_opt\_ml}
\date{2026-11-10}

\begin{document}
\maketitle

\section*{Summary}
The goal is to set up many CCDs for low noise ($\se$) and low dark charge ($\lambda$) with as few images per CCD as
possible. The first year asked machine learning to predict each image's labels from a trace, using hundreds of labelled
images from one module. That works for the noise and not for the dark charge, for reasons that are about the problem setup,
not the method (Section~\ref{sec:why}). This plan points the machine learning at the parts of the problem where it can
win:
\begin{enumerate}
\item \textbf{Learn across CCDs (the fleet).} A model trained on the CCDs already set up, which adapts to a new CCD from
  one or a few of its images (\emph{few-shot}). The ``one calibration image'' result of board 102 is the simplest case.
\item \textbf{Let the model choose what to measure next} (Bayesian optimisation, with the fleet model as prior): the
  standard machine-learning tool for ``each measurement is expensive, find the best setting in few of them''.
\item \textbf{Cheap measurements first}: traces (seconds) and level-1 strips (under a minute) carry most of the noise
  and, for strips, the dark charge; images (10 minutes) only to confirm.
\item \textbf{Physics for the obvious part, machine learning for the rest}: the timing effects are known physics; the
  model learns what is left (voltage effects, interactions, CCD-to-CCD differences).
\end{enumerate}
The success measure is one number per CCD: \textbf{images needed to come within a few \% of its best setting, on noise and
dark charge together}, and how that number falls as the fleet grows. Four tests (T1--T4) estimate it, and each part of
it, on the v2 data already taken; the script is \texttt{lvl1/fleet\_tests.py} (Section~\ref{sec:tests}).

\section{Why the first setup was hard}\label{sec:why}
What worked: the traces predict the noise of an amplifier never seen in training about twice as well as the settings
(9.7\,\% against 18.1\,\% on session E), and the direction of clear noise changes (88\,\% against 72\,\%).
What did not, and why:
\begin{itemize}
\item \textbf{The dark charge is not in the input.} A trace is a 0.19\,s reading of the output amplifier with almost no
  charge in it; the dark charge is electrons made in the CCD and its serial register during a 10-minute readout. No
  model can predict what its input never recorded. (A strip does record it: it counts electrons.)
\item \textbf{Too little data for the question.} 50--130 labelled packets, 3--4 amplifiers of one module, 15 knobs, with
  voltage effects of a few \% next to label noise of 2--10\,\%. ``A new amplifier'' was close to one trial.
\item \textbf{Simple physics already explains most of it.} A settings-only model explains 91\,\% of the dark-charge
  changes, almost all through \texttt{delayInt}; what is left is close to the label noise. There was nothing left for
  machine learning to find there.
\item \textbf{The prize near nominal is small}: on the first night the good settings were within about 3\,\% of each
  other in noise, while the models' error is about 13\,\%.
\end{itemize}
Machine learning pays off when (1) the input carries the information, (2) there is a lot of data relative to the
complexity, and (3) simple physics does not already give the answer. Commissioning many CCDs can meet all three: the
fleet supplies the data, the strips supply dark-charge information, and the model is pointed at the non-obvious part.

\section{One calibration image, and learning across the fleet}
Board 102 sat at nominal all night. Every model trained on board 101 predicted its levels badly ($\lambda$ about $-90$\,\%,
$\se$ physics model $-30$\,\%). After correcting each model by its miss on board 102's \emph{first} image, every model was
at the label scatter ($\se$ about 11\,\%, $\lambda$ about 8\,\%, gain 0.3\,\%). Each amplifier has its own level (an
offset), and one image measures it. (Caveat: 102 never changed settings, so this showed that the level is fixed by one
image, not that the response to the settings transfers.)

This is the simplest form of the first idea. Write each label of amplifier (later: CCD) $a$ at settings $\theta$ as
\[
\log y_a(\theta) = \underbrace{f(\theta)}_{\text{shared by the fleet}} + \underbrace{c_a}_{\text{level}} +
\underbrace{g_a(\theta)}_{\text{what is particular to } a}.
\]
The fleet teaches $f$, and how large $c_a$ and $g_a$ typically are. A new CCD needs only enough images to pin down its own
$c_a$ (one image) and the part of $g_a$ that matters. A \emph{multi-task Gaussian process} learns this split directly:
its kernel is the shared response times an amplifier similarity $\rho$ (1 = same shape), plus an amplifier-offset term,
and $\rho$ and the offset size are learned from the fleet. With more CCDs in the fleet, $f$ and the priors on $c_a$,
$g_a$ get better, and a new CCD needs fewer images. That is the property worth demonstrating: \textbf{images per CCD fall
as the fleet grows}.

\section{The obvious and the non-obvious parts}
\textbf{Obvious}: effects that are large, smooth, the same on every amplifier and predictable from physics. The timing is
the example: longer integration lowers the white noise and raises the slow noise; longer readout makes more dark charge
($\lambda \propto \texttt{delayInt}^{0.5}$ on every amplifier in v2). A settings model or a formula gets these from a handful
of points; machine learning adds nothing, as you said.

\textbf{Non-obvious}: what you cannot write down in advance, and what differs between CCDs:
\begin{itemize}
\item small voltage effects of a few \% each, which add up over 15 knobs and interact (in v2 the corner
  \texttt{delayRGH} high + VSUB high had $\lambda$ 80\,\% above the settings model on every amplifier: an interaction no
  one predicted);
\item effects that differ between CCDs (VSUB's effect on $\lambda$ differed 8$\times$ between amplifiers in v2: strong on
  ch1 and ch3, flat on ch0 and ch2);
\item whatever is left after the physics prediction.
\end{itemize}
In practice: give the model the physics as a baseline (or as its mean function) and let it learn the residual, and learn
from the fleet \emph{which knobs are CCD-specific}: those are the only ones a new CCD needs measured. Test T4 measures
this split on v2.

\section{The plan: four parts}
\begin{description}
\item[M1 Fleet model (few-shot).] Multi-task GP (and simpler references: shared GP + offset, trace trees + offset) trained
  on all amplifiers/CCDs; adapts to a new one from $k$ of its images. Grows with every CCD commissioned.
\item[M2 Sequential search.] Bayesian optimisation with M1 as prior: after each image, choose the next setting by
  expected improvement on $J(w) = w \log\se + (1-w)\log\lambda$. Start from one image at nominal (the calibration image).
  Replaces the funnel: it ranks and updates as images come in, instead of ranking once.
\item[M3 Cheap measurements.] Traces of every candidate setting (seconds each) as extra inputs to M1/M2; level-1 strips
  as labels if they track the image labels (T3). The search then spends minutes per evaluation, not ten.
\item[M4 Physics + residual.] The timing physics as the baseline; machine learning on the residual; the list of
  CCD-specific knobs from the fleet.
\end{description}
Optional, cheap: trace-based checks for broken amplifiers before any image (board 102 ch0/ch1 were dead; the traces
show it).

\section{Tests on the data already taken}\label{sec:tests}
All four tests treat one amplifier at a time as the ``new CCD'' and the other amplifiers as the fleet. \textbf{The new
amplifier is scored on one block of packets in time, and the fleet trains only on the other blocks}: the amplifiers of a
packet share its moment (drift, a stalled queue), which a CCD measured on another day would not share. (Without this,
the v1 run below showed $\lambda$ on a ``new'' amplifier as easier than on a trained one: a leak.) Script:
\texttt{lvl1/fleet\_tests.py}; driver: \texttt{lvl1/run\_fleet\_tests.sh}.

\subsection*{T1 Few-shot: how many images of a new CCD does the model need?}
Leave one amplifier out; give the model $k = 0, 1, 2, 4, 8$ of its images (random, from other time blocks); score its
predictions on one block. Output: error of $\se$, $\lambda$, gain against $k$; the regret of the setting each model would
pick (how much worse than the best in the block), for $\se$, $\lambda$ and $J(0.5)$; the multi-task GP's learned
$\rho$ and offset size. \textbf{Pass (proposed):} at $k = 1$ the $\se$ error within $2\times$ the nominal label noise and
the $\lambda$ error within $2\times$ its label noise; regret below 3\,\% ($\se$) and 10\,\% ($\lambda$) by $k \le 4$.

\subsection*{T2 Bayesian optimisation replayed on the measured packets}
For the left-out amplifier, its measured settings in one time block are the candidates (about 30 per block in v2). A
strategy picks the next setting, ``takes the image'' (reveals the label), updates, and so on. Strategies: random order
(200 draws); the fleet model's ranking without updates; BO with the fleet as prior; BO with the fleet plus the traces of
every candidate; BO with no fleet data. Output: best-so-far against the number of images; images needed to come within
2\,\% and 5\,\% of the best in the pool, for $w = 1, 0.5, 0$. \textbf{Pass (proposed):} BO with the fleet reaches 2\,\%
in at most a third of the images random order needs, and in at most 10. Limits: the pool is the settings v2 happened to
measure, so ``best in pool'' is not the true optimum, and a pool of 30 caps the gain.

\subsection*{T3 Can the strip stand in for the image?}
Per target, on the same packets: strip against image (correlation; error after an amplifier offset; error of a
calibration model $\log y_\text{image} \sim \log y_\text{strip} + \text{amplifier}$, cross-validated by packet);
repeatability at nominal of each; direction of the within-packet strip change (new against reference strip) against the
image change, on clear changes. \textbf{Pass (proposed):} calibrated error at or below the image's own repeatability,
change direction $\ge 90$\,\%. If $\lambda$ passes, the search can run on strips (about 1 minute) and image only to
confirm.

\subsection*{T4 What is obvious and what is left}
Amplifier offsets removed (one calibration image), then: variance explained by the timing knobs alone; by all settings;
what the voltages explain of the timing residual on trained amplifiers and on a new one; what is left against the label
noise; and per voltage knob, the slope all amplifiers share against the amplifier-to-amplifier spread (flagged
amplifier-specific when the spread exceeds twice the slope error). \textbf{Decision:} knobs flagged amplifier-specific
are the ones a new CCD must measure; if the voltages explain nothing of the $\lambda$ residual on a new amplifier, the
dark-charge search is timing + those few knobs.

\subsection*{What the v1 run shows (code check, not the evidence)}
The script runs on v1 (43 packets, 4 amplifiers, one knob at a time, 20 nominal repeats per amplifier); about 90\,s. v1 is
the wrong design for these questions (most voltages moved by one max\_step, pools of about 20 candidates per block), so
these numbers check the code, not the plan:
\begin{itemize}
\item T1: one image of the new amplifier takes the shared GP from 13.3 to 8.9\,\% on $\se$ and from 45 to 26.5\,\% on
  $\lambda$; more images add little (8.2 and 23.2\,\% at $k = 8$). The multi-task GP learns $\rho = 1$: on v1 the
  amplifiers share the response shape and differ in level, which is the structure that makes one image enough.
\item T2: inconclusive on v1 (small pools, many nominal repeats); BO with fleet + traces reached 2\,\% of the best
  $\lambda$ setting in 5 images against 13 for random order.
\item T4: after the offsets, timing explains 66\,\% of the $\se$ variance; the voltages add nothing on v1 (they barely
  moved); $\lambda$ is barely explained by any setting (26\,\% left against 10\,\% label noise): on the first night most
  of its variation was the CCD's state, not the settings.
\end{itemize}

\section{After the tests}
\begin{itemize}
\item \textbf{T1 and T2 pass on v2:} the decisive test is live: run M2 on a board the fleet has not seen (board 102 with
  working amplifiers, or the next CCD), count images to the best setting, and compare with your usual procedure. More
  CCDs then help more than more packets on one CCD.
\item \textbf{T3 passes for $\lambda$:} make the strip the default label of the search.
\item \textbf{T4:} the CCD-specific knobs set the search space per CCD; everything else is set once per CCD type.
\item \textbf{Data still worth taking} (lvl1/packet\_plan\_pre): the settling and exposure pre-tests, because they decide
  whether the labels the search optimises mean what we think. The large v3 sets are only needed if T1--T4 show the
  voltage response is the bottleneck.
\end{itemize}

\section{Commands}
\begin{Verbatim}
# on the machine with the v2 processed data (dataset.csv), from the repository root:
DATA=processed_v2 bash lvl1/run_fleet_tests.sh          # all runs, ~5-20 min; results in processed_v2/lvl1_results/fleet_tests_<date>/
# single tests:
.venv/bin/python lvl1/fleet_tests.py fewshot --data processed_v2 --plan lvl1/packet_plan_v2/plan.json --sessions D E --channels 0 1 3 --drift --out <dir>
.venv/bin/python lvl1/fleet_tests.py bo       ... ;  ... strips ... ;  ... residual ...
# the v1 check:
.venv/bin/python lvl1/fleet_tests.py all --data processed --plan lvl1/packet_plan/plan.json --sessions A B C --channels 0 1 2 3 --drift --out <dir>
\end{Verbatim}
Send back the output folder; the \texttt{.md} files hold the tables.

\end{document}
